\documentclass[10pt,a4paper]{amsart}
\usepackage[margin=19mm]{geometry}
\usepackage{amsmath,amssymb,mathtools}
\usepackage[scr=rsfs,cal=euler]{mathalfa}
\usepackage{xfrac}
\usepackage[unicode,hidelinks,bookmarks=false]{hyperref}
\newcommand{\ie}{i.e.\ }
\newcommand{\cf}{cf.\ }
\newcommand{\resp}{respectively\ }
\newcommand{\Subsection}{\subsection}
\providecommand{\nobreakdash}{\nobreak\hbox{-}}
\setlength{\emergencystretch}{2em}
\multlinegap0pt
\usepackage{amssymb}
\usepackage{tikz}
\usepackage{algorithm}\usepackage{algpseudocode}
\usepackage{xcolor}
\usepackage{mathtools}
\usepackage[shortlabels]{enumitem}
\newtheorem{proposition}{Proposition}
\title{Exact curve counting of given word length \\ on the once-punctured torus}
\author{Filippo Baroni, David Fisac, Mingkun Liu}
\date{Source: arXiv:2609.01382v1 (CC BY 4.0)}
\begin{document}
\maketitle

\noindent\emph{Source and licence.} Exact curve counting of given word length \\ on the once-punctured torus. Version arXiv:2609.01382v1 (CC BY 4.0), distributed by arXiv under the Creative Commons Attribution 4.0 International licence, http://creativecommons.org/licenses/by/4.0/, as recorded in the arXiv licence metadata for this exact version and shown on the abstract page https://arxiv.org/abs/2609.01382v1. Full licence notice: \texttt{licenses/arxiv-2609.01382-CC-BY-4.0.txt}.\par\smallskip

\noindent\textit{Editorial scope.} Counted unit: the article's proposition proposition (source0.tex lines 1151--1159). The article's complete original proof of it is delivered with it (source0.tex lines 1160--1187, one \texttt{proof} environment of 28 lines). No mathematics is written, repaired or re-derived: the statement and the proof are the article's own lines, in the article's own notation, and no retained line is substituted or re-flowed.  No further source-local context is needed: every cross-reference of every retained line resolves inside the two delivered spans.  Nothing was omitted from any retained span, so the package carries no omission record.  No retained line contains a citation, so the wrapper carries no bibliography and no bibliography entry.  No retained line contains a figure environment, an image-inclusion command or drawing code, so no image is delivered and the package declares no asset dependency.  Editorial formatting only, in addition: an AMS \texttt{amsart} preamble that restates the article's own declarations and macros used by the retained lines, namely the environments \texttt{proposition} on separate counters, together with the standard packages those lines need.  The retained environments print in this document as Proposition 1 in order of appearance, whereas the article prints the same items with the numbering of its own full document.  The complete native source of the article as distributed in the arXiv e-print for this exact version is bundled unchanged as \texttt{sources/209106/source0.tex}.
	\begin{proposition}
		Let $w$ be a cyclically reduced cyclic word in $\{\mathtt{a},\mathtt{b},\mathtt{z},\mathtt{A},\mathtt{B},\mathtt{Z}\}$.
		If $w$ is allowed, then the following claims hold.
		\begin{enumerate}
			\item If $w$ admits a full commutator move $m$, then $m(w)$ is allowed;
			\item If $w$ admits no full commutator moves but admits a partial commutator move $m$, then $m(w)$ is allowed;
			\item If $w$ is a branch word, then $L^*(w)$ and $R^*(w)$ are allowed.
		\end{enumerate}
	\end{proposition}

	\begin{proof}\leavevmode We will prove it point by point.
		\begin{enumerate}
			\item Let us do only (F1), since the other full commutator moves are symmetric.
			Let $u=m(w)$.
			Write $w=\mathtt{aB}^k\mathtt{A b} v$ for some reduced word $v$, so that $u=(\mathtt{zB})^{k-1}\mathtt{z}v$.
			Since $w$ is allowed, $v$ does not start or end with $\mathtt{Z}$, so $u$ has no $\mathtt{zZ}$ or $\mathtt{Z z}$ subwords.
			Since $w$ is cyclically reduced, it does not start with $\mathtt{B}$ or end with $\mathtt{A}$.
			In particular, $u$ does not have $\mathtt{zB a}$ or $\mathtt{bA z}$ subwords (since $v$ also does not).
			We deduce that $u$ is allowed.
			\item Let us do only (P1), since the other partial commutator moves are symmetric.
			Let $u=m(w)$.
			Write $w=\mathtt{aB}^k\mathtt{A} v$ for some reduced word $v$, so that $u=(\mathtt{zB})^k v$.
			Because $w$ admits no commutator moves and is reduced, we know that $v$ does not start with $\mathtt{a}$ or $\mathtt{b}$, and it does not end with $\mathtt{Z}$ or $\mathtt{A}$. We check then that none of the forbidden subwords are in $u$.
			
			For $\mathtt{zZ}$, and $\mathtt{Zz}$: the block $(\mathtt{zB})^k$ contains no $\mathtt{Z}$ and $v$ does not begin with $\mathtt{Z}$. For $\mathtt{zBA}$: within $(\mathtt{zB})^k$ each $\mathtt{z}$ is followed by $\mathtt{B}$ and then $\mathtt{z}$, never by $\mathtt{a}$, and at the junction it is excluded because $v$ cannot start with $\mathtt{a}$. For $\mathtt{bAz}$: $u$ contains a $\mathtt{z}$ either inside $(\mathtt{zB})^k$ or inside $v$, since the letter before each $(\mathtt{z}B)$-block is either $B$ or it is the boundary, in which case it cannot be $\mathtt{bA}$ because $w$ is cyclically reduced. For $\mathtt{ZaB}$ and $\mathtt{AbZ}$: these involve $\mathtt{Z}$, which happens only inside $v$, and then it is controlled by $w$ being allowed.
			
			
			\item Let us do only $R^*$, since $L^*$ is symmetric.
			Let $u=R^*(w)$.
			This is again a case check using the table.
			It is clear that $u$ has no $\mathtt{zZ}$ or $\mathtt{Z z}$ subwords.
			A $\mathtt{zB a}$ subword of $u$ can only come from a $\mathtt{zB a}$ of $w$ (forbidden).
			A $\mathtt{Z aB}$ subword can only come from $\mathtt{Z aB}$ (forbidden).
			A $\mathtt{bA z}$ subword can only come from $\mathtt{bA z}$ (forbidden).
			An $\mathtt{A bZ}$ subword can only come from $\mathtt{A bZ}$ (forbidden).
			So $u$ is allowed. \qedhere
		\end{enumerate}
	\end{proof}

\end{document}
